%! Carrying Out a Test for the Difference of Two Population Proportions — Unit 6, Lesson 6.11 — Homework
\documentclass[10pt]{article}
\usepackage{apstats-article}
\usepackage{apstats-boxes}
\usepackage{grading}
\gradingassignment{apstats-u06-l11-hw}
\gradingcourse{advanced_statistics}
\newcommand{\TallMath}[1]{$\displaystyle #1\rule[-1.4em]{0pt}{3.2em}$}

\begin{document}

\pageheader{Unit 6, Lesson 6.11}{Homework}

\vspace{0.15in}

Complete the full four-step procedure for each problem. Use $\alpha = 0.05$ unless stated otherwise.

\begin{enumerate}

\item \textbf{Social Media and Academic Performance.} A researcher hypothesizes that
students who use social media more than 2 hours per day have a higher proportion of
low-GPA outcomes than those who use it 2 hours or less per day.
High use: 54 of 180 had a low GPA.
Low use: 36 of 200 had a low GPA.

\begin{enumerate}[label=\alph*.]
  \item State $H_0$ and $H_a$ (one-sided). \gradedhide{\writeline\\\gradedhere{\gradedblank[id=hw1a, type=rubric, points=2, display={$p_H, p_L$: true low-GPA proportions (high / low use). $H_0\colon p_H = p_L$; $H_a\colon p_H > p_L$.}]{\linewidth+6pt}{\gradedlinesheight{1}}}\writelines{9}}
  \item Compute $\hat p_c$ and verify conditions. \gradedhide{\writeline\gradedright{$\hat p_c = $~\gradedblank[id=hw1b1, type=short, points=1, answer=0.237, check=numeric, tolerance=0.002]{0.8in}{0.17in}}\\\gradedhere{\gradedblank[id=hw1b2, type=rubric, points=3, display={Independent random samples assumed; with $\hat p_c = 0.237$: $42.7, 137.3, 47.4, 152.6$ all $\ge 10$; samples $\le 10\%$ of all students.}]{\linewidth+6pt}{\gradedlinesheight{2}}}\writelines{9}}
  \item Compute $z$ and the $p$-value. \writelines{2}\\\gradedhide{\vspace{1pt}\noindent\textcolor{linegray}{\hrulefill}\gradedright{$z = $~\gradedblank[id=hw1c1, type=short, points=2, answer=2.74, check=numeric, tolerance=0.02]{0.8in}{0.17in}\quad $p$-value $=$~\gradedblank[id=hw1c2, type=short, points=2, answer=0.003, check=numeric, tolerance=0.001]{0.8in}{0.17in}}\vspace{3pt}}
  \newpage
  \item Conclude in context. \gradedhide{\writeline\\\gradedhere{\gradedblank[id=hw1d, type=rubric, points=2, display={$0.003 \le 0.05$: reject $H_0$; convincing evidence heavy (over 2 h) users have a higher low-GPA proportion.}]{\linewidth+6pt}{\gradedlinesheight{1}}}\writelines{9}}
\end{enumerate}

\item \textbf{Employee Turnover.} A company compares turnover rates for employees who
went through a new onboarding program versus the old one.
New program: 18 of 240 left within the first year.
Old program: 30 of 240 left within the first year.
Test whether the new program reduced turnover (one-sided, $\alpha=0.05$).

\begin{enumerate}[label=\alph*.]
  \item State $H_0$ and $H_a$. \gradedhide{\writeline\\\gradedhere{\gradedblank[id=hw2a, type=rubric, points=2, display={$p_N, p_O$: true first-year turnover proportions. $H_0\colon p_N = p_O$; $H_a\colon p_N < p_O$.}]{\linewidth+6pt}{\gradedlinesheight{1}}}\writelines{9}}
  \item Compute $\hat p_c$ and verify conditions. \gradedhide{\writeline\gradedright{$\hat p_c = $~\gradedblank[id=hw2b1, type=short, points=1, answer=0.100, check=numeric, tolerance=0.002]{0.8in}{0.17in}}\\\gradedhere{\gradedblank[id=hw2b2, type=rubric, points=3, display={Random samples assumed; with $\hat p_c = 0.1$: $24, 216$ (each group) all $\ge 10$; samples $\le 10\%$ of all employees.}]{\linewidth+6pt}{\gradedlinesheight{2}}}\writelines{9}}
  \item Compute $z$ and the $p$-value. \writelines{2}\\\gradedhide{\vspace{1pt}\noindent\textcolor{linegray}{\hrulefill}\gradedright{$z = $~\gradedblank[id=hw2c1, type=short, points=2, answer=-1.83, check=numeric, tolerance=0.02, display={\boldmath$-1.83$}]{0.8in}{0.17in}\quad $p$-value $=$~\gradedblank[id=hw2c2, type=short, points=2, answer=0.034, check=numeric, tolerance=0.002]{0.8in}{0.17in}}\vspace{3pt}}
  \newpage
  \item Conclude in context. \gradedhide{\writeline\\\gradedhere{\gradedblank[id=hw2d, type=rubric, points=2, display={$0.034 \le 0.05$: reject $H_0$; convincing evidence the new program has lower first-year turnover.}]{\linewidth+6pt}{\gradedlinesheight{1}}}\writelines{9}}
\end{enumerate}

\end{enumerate}

\end{document}
